\section{Discussion}
\label{sec:discussion}

\begin{table}[t]
    \centering
    \resizebox{\textwidth}{!}{%
    \begin{tabular}{@{}p{6.2cm}p{4.6cm}p{6.2cm}@{}}
        \toprule
        \textbf{Claim} & \textbf{Verdict} & \textbf{Evidence} \\
        \midrule
        Adaptive protocols identify more vulnerabilities than static ones at matched budget
          & \textbf{Supported, strongly}
          & E1: $0.00 \to 1.00$; E2: $+0.65$ on \gptfour; E4 (live): $+0.72$ on \qwenseven \\
        \addlinespace
        The mechanism is using nulls as iterative triggers for redesign
          & \textbf{Refuted within a search; supported across a portfolio}
          & E1: $\Delta = +0.02$ to $-0.12$ vs a blind schedule; E1c: AUC $\approx 0.5$. E2 and E4:
            Thompson beats round-robin decisively \\
        \addlinespace
        Adaptive protocols provide better evidence for genuine robustness
          & \textbf{Refuted as currently practised}
          & E3: 63--67\% coverage of a nominally 95\% bound; anytime-validity does not fix it.
            E4 shows the design-space ceiling is the honest substitute \\
        \bottomrule
    \end{tabular}
    }
    \caption{Verdict on the three claims contained in the hypothesis. The hypothesis is right about
    the conclusion and wrong about the mechanism, and the correction is practically consequential.}
    \label{tab:verdict}
\end{table}


\subsection{What the results mean}

The hypothesis we set out to test contains three claims, and they receive different verdicts
(\tabref{tab:verdict}). Adaptivity does find substantially more at matched compute, and the effect
is as large as anyone could want. But almost none of it comes from the mechanism the hypothesis
names. What matters is searching over \families at all, and learning across a portfolio which
family works. Using the continuous signal inside a failing run as a redesign trigger is worth
nothing and can be actively harmful.

The correction is practically consequential rather than semantic. A team that instruments its
existing single-design pipeline with a stall detector will get nothing, and at large budgets will
lose detections. A team that reallocates the same budget across three designs will go from 0\% to
100\%.

\subsection{Why the two levels of adaptivity differ}

{\bf Why does the same idea work at one level and fail at the other?} Because they face different
signal-to-noise regimes. At the portfolio level, the evidence accumulating against a method is a
run of independent Bernoulli nulls across \emph{different} behaviours. Twenty nulls in a row from a
method is overwhelming evidence that the method does not work on this model, and Thompson sampling
exploits that correctly. At the within-search level, the ``signal'' is a single monotone surrogate
on one behaviour, and E1c shows it does not separate the surviving-null population. The lesson
generalises beyond our substrates: \emph{nulls are informative in aggregate across independent
trials, and uninformative as a within-trial progress signal.}

This also reconciles our result with \citet{andriushchenko2024adaptive}, whose $0\% \to 50\% \to
100\%$ progression on \llamaseven is often read as evidence for adaptive \emph{search}. Our
decomposition says the gain is a \family effect: self-transfer initialisation and template
switching are changes of design, not data-driven adjustments within one design. Their own R2D2
negative control agrees, since more self-transfer gives 12\% there while a template-family switch
gives 90\%.

\subsection{Failure modes}

\para{When portfolio adaptivity loses.} On \llamaseven and Vicuna-13B a dominant method exists, so
exploration is wasted and blind round-robin wins by $0.185$. The practical diagnostic is simple: if
you expect a wide and unknown method-quality spread, adapt; if you suspect a dominant method,
do not.

\para{When the ladder loses.} The ``persist while the surrogate is warm'' rule keeps budget on rungs
whose surrogate is high but whose judged outcome is negative. Because the surrogate is
non-diagnostic on the surviving-null population, this is a pure loss, and it grows with budget.

\para{When the static protocol is actively misleading.} E4 shows the worst case is not
under-powering but mis-estimation of an intervention's \emph{effect size}. Judged by P1 the
anti-sycophancy intervention on \qwenseven looks nearly inert ($0.077 \to 0.037$); judged by P4 it
looks strongly effective ($1.000 \to 0.519$). Single-family evaluation gets the direction of the
conclusion wrong, not just the precision.

\para{Judge disagreement is large and is not noise.} The configuration
\texttt{llama2-7b\_icl\_one\_shot} scores $0.66$ under the GPT-4 judge and $0.10$ under the
rule-based judge. Consistent with the theoretical ceiling on judge substitution~\cite{dorner2024limits},
any single-judge efficiency claim in this literature should be treated as provisional. We report
both outcomes throughout.

\subsection{The surprise}

We expected anytime-valid inference to substantially repair coverage in E3, and it barely moved it
($0.638 \to 0.671$). The problem in safety evaluations is not primarily that people peek at their
data. It is that the estimand they bound, the rate observed at budget $t$, is not the estimand they
report, elicitability. No amount of sequential rigour fixes an estimand error. This suggests the
useful research target is a bound whose estimand is elicitability \emph{under a specified design
space}, evaluated for coverage against held-out \families.

\subsection{Limitations}

We state these plainly because several of them bound the conclusions materially.

\begin{itemize}[leftmargin=*,itemsep=1pt,topsep=2pt]
    \item \textbf{Retrospective replay.} E1--E3 replay logs generated by other people's design
    choices, so a \family absent from those logs cannot be evaluated. \staticzero's total failure
    is partly a property of the specific \texttt{plain\_init} ablation, which is harsher than the
    headline ``Prompt + RS'' row reporting 50\% in the original paper. We do \emph{not} claim to
    reproduce that table.
    \item \textbf{Conservative replay semantics.} A run the original authors ended without success
    is treated as a failure even if budget remained, which caps the persistence arms. This is
    applied identically to every arm, so it does not favour the adaptive ones; if anything it makes
    \fixedsplit and \ladder look more similar than they are.
    \item \textbf{Unpaired rung 2.} Goal strings are unrecoverable for the in-context-template
    configurations, so the third rung is matched by resampling, which is an exchangeability
    assumption. The exactly-paired two-rung variant gives the same qualitative answer
    ($\Delta = +0.98$ vs \staticzero).
    \item \textbf{One realisation per cell} in the JailbreakBench matrix, so all E2 uncertainty is
    policy-randomisation and bootstrap, not attack re-run variance.
    \item \textbf{Heuristic run boundaries} in the trajectory parse, accurate to within $\pm 2$ of
    the expected 50 runs per configuration.
    \item \textbf{The surrogate is an approximation.} E1c uses target-token probability, which
    \citet{jones2025forecasting} themselves use for most experiments, but the mapping to elicitation
    probability is approximate. A stronger signal might be diagnostic where this one is not. This is
    the most important threat to the E1c conclusion and therefore to the E1 verdict.
    \item \textbf{E3's ground truth} is the end of the observed trace, not true elicitability at
    infinite budget, so the reported coverage is if anything optimistic.
    \item \textbf{No frontier API access}, so the live arm uses local open-weights models. Results
    may not transfer to frontier models with different safety training.
    \item \textbf{E4 is small}: 32 items per model, 19--27 eligible after the correctness filter,
    and four pressure families. The family-level contrasts are large enough to survive this
    ($0.000$ vs $0.852$), but the per-family rates carry wide intervals, and the ``ceiling'' is a
    ceiling over the four families we wrote, not over all possible designs. That is exactly the
    caveat this paper argues every null should carry.
    \item \textbf{\gsmk contamination} cannot be excluded. It is mitigated by measuring answer
    \emph{flipping under pressure} rather than accuracy, and by conditioning on items the model
    already answers correctly, but it is not eliminated.
\end{itemize}

\subsection{Broader impact}

This work is defensive evaluation methodology. E1--E3 replay already-published benchmark artifacts
and generate no new attack content, and E4 uses a non-harmful failure mode, sycophantic
answer-flipping on grade-school arithmetic, and produces no jailbreak artifacts. The main risk we
see is the opposite of misuse: our headline finding could be read as a licence to distrust all null
results. That is not the claim. The claim is that a null should be reported together with the design
space that was searched, and that a binomial interval computed on a null does not bound a model's
vulnerability. Used that way, the finding should make safety cases more honest rather than less
usable.
